% SOURCE_AUTHORITY: sga5_fr_workpass.tex, lines 2745--2893 (LF-normalized unit).
% SOURCE_SHA256: 791F4EFFC5E02832D5D77ED03518C8156D6F07E4C8238B03545DB93D883FBB28
\subsection*{4.1}
Sejam, para $i=1,2$, $X_i$ um $S$-esquema, $X=X_1\times_SX_2$, $p_i:X\to X_i$ as projeções, $L_i\in\ob\Dctf(X_i)$ e $K_{X_i}$ o complexo dualizante (1.3). As flechas canônicas
\begin{equation}
p_i^*\uRHom(L_i,K_{X_i})\overset{\mathbf L}{\otimes}p_i^*L_i
\longrightarrow p_i^*K_{X_i},
\tag{4.1.1}
\end{equation}
compostas de 2.1.1 e da flecha de avaliação 2.2.2, fornecem, tomando o produto tensorial sobre $X$ e compondo com o isomorfismo de Künneth 1.7.6, um emparelhamento
\begin{equation}
(DL_1\overset{\mathbf L}{\otimes}_S L_2)\overset{\mathbf L}{\otimes}(L_1\overset{\mathbf L}{\otimes}_S DL_2)
\longrightarrow K_X,
\tag{4.1.2}
\end{equation}
que pode ser reescrita, graças a 3.1.1, 3.1.2,
\begin{equation}
\uRHom_S(L_1,L_2)\overset{\mathbf L}{\otimes}\uRHom_S(L_2,L_1)
\longrightarrow K_X.
\tag{4.1.3}
\end{equation}
Daí se deduz um emparelhamento, denominado produto cup,
\begin{equation}
\langle\ ,\ \rangle:\Hom_S(L_1,L_2)\otimes\Hom_S(L_2,L_1)
\longrightarrow \H^0(X,K_X).
\tag{4.1.4}
\end{equation}
Também se pode considerar 4.1.2 como o produto tensorial externo das flechas de avaliação
\[
DL_i\overset{\mathbf L}{\otimes}L_i\longrightarrow K_{X_i}.
\]
Levando em conta o isomorfismo de Künneth 2.2.4, vê-se assim que 4.1.2 é uma ``dualidade perfeita'', isto é, identifica, mediante o isomorfismo de adjunção 2.2.1, cada um dos complexos $DL_1\overset{\mathbf L}{\otimes}_S L_2$, $L_1\overset{\mathbf L}{\otimes}_S DL_2$ com o dual do outro (com valores em $K_X$).

Quando $X_1=X_2=S$, 4.1.3 é um caso particular de um emparelhamento válido para complexos perfeitos sobre um topos anelado qualquer (SGA 6 I 7.8), e, com as notações de (SGA 6 I 8.3), tem-se
\begin{equation}
\langle u,v\rangle=\mathrm{Tr}(vu)=\mathrm{Tr}(uv).
\tag{4.1.5}
\end{equation}

\subsection*{4.2}
Sejam $c:C\to X$, $d:D\to X$ e $e=c\times_Xd:E=C\times_XD\to X$. Para $P,Q\in\ob\Dctf(X)$, dispõe-se de uma flecha canônica (que não é um isomorfismo em geral)
\begin{equation}
c^!P\overset{\mathbf L}{\otimes}_X d^!Q
\longrightarrow e^!(P\overset{\mathbf L}{\otimes}Q),
\tag{4.2.1}
\end{equation}
definida, do mesmo modo que 1.7.3, como adjunta da flecha composta
\[
e_!(c^!P\overset{\mathbf L}{\otimes}_X d^!Q)
\xrightarrow[(1)]{\sim}
 c_!c^!P\overset{\mathbf L}{\otimes}d_!d^!Q
\xrightarrow{(2)} P\overset{\mathbf L}{\otimes}Q,
\]
onde (1) é o isomorfismo de Künneth (SGA 4 XVII 5.4.2.2) e (2) o produto tensorial das flechas de adjunção. Aplicando $e_!$ a 4.2.1 e compondo com o isomorfismo de Künneth que acabamos de citar, obtém-se uma flecha canônica
\begin{equation}
c_!c^!P\overset{\mathbf L}{\otimes}d_!d^!Q
\longrightarrow e_!e^!(P\overset{\mathbf L}{\otimes}Q).
\tag{4.2.2}
\end{equation}
Da mesma forma, compondo a flecha de Künneth
\[
c_*c^!P\overset{\mathbf L}{\otimes}d_*d^!Q
\longrightarrow e_*(c^!P\overset{\mathbf L}{\otimes}_X d^!Q)
\]
com a flecha derivada de 4.2.1 aplicando $e_*$, obtém-se uma flecha canônica
\begin{equation}
c_*c^!P\overset{\mathbf L}{\otimes}d_*d^!Q
\longrightarrow e_*e^!(P\overset{\mathbf L}{\otimes}Q).
\tag{4.2.3}
\end{equation}
Para $P=DL_1\overset{\mathbf L}{\otimes}_S L_2$, $Q=L_1\overset{\mathbf L}{\otimes}_S DL_2$, de 4.2.2, 4.2.3 deduzem-se, por composição com 4.1.2 e com o isomorfismo $e^!K_X=K_E$, e com as identificações 3.1.1, 3.2.1, os emparelhamentos
\begin{equation}
c_!\uRHom(c_1^*L_1,c_2^!L_2)\overset{\mathbf L}{\otimes}
 d_!\uRHom(d_2^*L_2,d_1^!L_1)
\longrightarrow e_!K_E,
\tag{4.2.4}
\end{equation}
\begin{equation}
c_*\uRHom(c_1^*L_1,c_2^!L_2)\overset{\mathbf L}{\otimes}
 d_*\uRHom(d_2^*L_2,d_1^!L_1)
\longrightarrow e_*K_E,
\tag{4.2.4'}
\end{equation}
\begin{equation}
\langle\ ,\ \rangle:
\Hom(c_1^*L_1,c_2^!L_2)\otimes\Hom(d_2^*L_2,d_1^!L_1)
\longrightarrow \H^0(E,K_E).
\tag{4.2.5}
\end{equation}
Esses emparelhamentos são compatíveis com a localização étale: dado um quadrado comutativo (não necessariamente cartesiano)
\[
\begin{tikzcd}[column sep=2.5em,row sep=1.5em]
& E' \arrow[dl] \arrow[dr] & && & E \arrow[dl] \arrow[dr] & \\
C' \arrow[dr] && D' \arrow[dl] & \text{ étale sobre } & C \arrow[dr] && D \arrow[dl] \\
& X &&&& X &
\end{tikzcd}
\]
se denotarmos por $c'$, $d'$, $e'$ as compostas $C'\to C\to X$, $D'\to D\to X$, $E'\to E\to X$, tem-se um quadrado comutativo.


\begin{equation}
\tag{4.2.6}
\begin{tikzcd}[column sep=large,row sep=large]
\Hom(c_1^*L_1,c_2^!L_2)\otimes\Hom(d_2^*L_2,d_1^!L_1)
  \arrow[r,"4.2.5"] \arrow[d] &
\H^0(E,K_E)\arrow[d] \\
\Hom(c_1'{}^*L_1,c_2'{}^!L_2)\otimes\Hom(d_2'{}^*L_2,d_1'{}^!L_1)
  \arrow[r] &
\H^0(E',K_{E'})
\end{tikzcd}
\end{equation}
onde as flechas verticais são as restrições, e a flecha horizontal inferior é a composta da flecha análoga a 4.2.5 e da restrição
\[
\H^0(E'',K_{E''})\longrightarrow \H^0(E',K_{E'}),
\]
com $E''=C'\times_{X}D'$ (têm-se quadrados comutativos análogos com flechas do tipo 4.2.4, 4.2.4', cuja escrita deixamos ao leitor). Se $Z$ é étale sobre $E$, para
\[
u\in\Hom(c_1^*L_1,c_2^!L_2),\qquad
v\in\Hom(d_2^*L_2,d_1^!L_1),
\]
escreveremos
\begin{equation}
\tag{4.2.7}
\langle u,v\rangle_Z\in \H^0(Z,K_Z)
\end{equation}
para a imagem de $\langle u,v\rangle$ em $\H^0(Z,K_Z)$ por restrição (na prática, $Z$ será uma componente conexa de $E$). Se $z$ é um ponto fechado isolado de $E$, $\bar z$ um ponto geométrico sobre $z$, de 4.2.6 deduz-se que
\[
\langle u,v\rangle_z\in \H^0(z,K_z)=A
\]
depende unicamente da situação obtida mediante localização estrita nas imagens de $\bar z$ em $C,D,X$.

\begin{statement}{Exemplo 4.3}
Suponhamos que $X_i$ seja suave, de dimensão relativa pura $r_i$, e que $c$ e $d$ sejam imersões fechadas, com $c_2$ (resp. $d_1$) finito. Para $L_1=A_{X_1}$, $L_2=A_{X_2}$, tem-se, por 3.1.3,
\[
DL_1\otimes_S^{\mathbf L}L_2=A_X(r_1)[2r_1],\qquad
L_1\otimes_S^{\mathbf L}DL_2=A_X(r_2)[2r_2],
\]
e 4.1.2 identifica-se com o produto
\[
A_X(r_1)[2r_1]\otimes^{\mathbf L}A_X(r_2)[2r_2]\longrightarrow A_X(r)[2r],
\]
onde $r=r_1+r_2$, e, portanto, tendo em conta 3.2.1, 4.2.5 identifica-se com o produto cup habitual
\[
\H_C^{2r_1}(X,A(r_1))\otimes \H_D^{2r_2}(X,A(r_2))
\longrightarrow
\H_{C\cap D}^{2r}(X,A(r)).
\]
Tomemos como correspondência cohomológica de $X_1$ a $X_2$ (resp. de $X_2$ a $X_1$) a classe de $c$ (resp. de $d$) (3.2 a)). De (SGA $4^{1/2}$ Ciclo 2.3.8) deduz-se que, em um ponto isolado $z$ de $C\cap D$, $\langle \cl(c),\cl(d)\rangle_z$ não é senão a imagem em $A$ da multiplicidade de interseção de $C$ e $D$ em $z$.
\end{statement}

Eis agora o resultado principal da exposição.
